\scanpage{43}
% A pale sweeping diagonal crosses the body of the sheet.
\tikz[remember picture,overlay]{\coordinate (scanXLIIIstart) at (1.2,0);}
le foncteur $Rf^!$ comme fourni par la \deleted{\unsure{correspondance|complexe}} donnée 3.1.10
\[
 U\longmapsto f_!^\bullet((\mathbb Z/n)_U)
\]
% slanted
\margin{$\otimes_{\mathbb Z/n}\mathcal A$ ? sinon, expliciter \unsure{de nouveau|à nouveau|le nouveau} $\mathcal A,\mathcal B$}
Rappelons\tikz[baseline]{\draw (0,-.06) -- (.18,.34); \draw (.04,.03) -- (.12,-.07);}\margin{réf ?} que la \uline{résolution plate canonique} \inserted{$L_\bullet(\underline F)$} d'un Module $\underline F$ s'obtient en représentant $\underline F$ comme quotient du Module engendré par le faisceau d'ensembles $\underline F$ : $\varepsilon:L_0(\underline F)\longrightarrow F$, en appliquant ensuite cette construction à $\mathrm{Ker}(\varepsilon)$, et en itérant le procédé.

Le foncteur $Rf^!$ est encore fourni, \inserted{par la méthode de 3.1.10,} par la donnée
\[
 U\longmapsto\inserted{K_0(U)\overset{\mathrm{dfn}}{=}}L_\bullet(f_!^\bullet(\mathbb Z/n)_U)\deleted{=K_0(U)}.
\]
Si maintenant $\underline G$ est un complexe de $\mathcal A$-Modules acycliques\unsure{,|.} $(\qquad\qquad)$, on a
\margin{??}
\[
 K_0^!(\underline G)\xrightarrow{\sim}RK_0^!(\underline G).
\]
On \deleted{en déduit}\inserted{retrouve} que le faisceau d'anneaux $\mathcal A$ ne joue qu'un rôle bidon dans la construction de $\deleted{\unsure{RK^!|RK_0^!}} Rf^!$, \deleted{on a}
\tikz[remember picture,overlay]{\draw[gray,thin] (scanXLIIIstart) to[bend left=4] (0,0);}
% archive sheet number: 42
